\section{The product-formula contradiction}\label{sec:contradiction}
For a nonzero rational $q$, unique factorization gives
$q=\pm\prod_pp^{v_p(q)}$, with only finitely many nonzero exponents.
Thus
\begin{equation}\label{eq:product-formula}
 \log|q|=\sum_pv_p(q)\log p.
\end{equation}
This is the arithmetic obstruction to a nonzero rational expression being
too small relative to its denominators and divisibility. The same general
principle underlies rational-approximation proofs such as
\cite{Apery,Beukers}; here it is applied to a selected determinant.

Assume all four target zeta values are rational. The rank argument selects
$D_n\in\Q\setminus\{0\}$, and the corrected-coordinate identification
makes it exactly the minor estimated analytically. For sufficiently large
$n\ge1$, the arithmetic estimate~\eqref{eq:P-certified} and the uniform
analytic estimate~\eqref{eq:uniform-upper} imply
\[
 P_*(100n)^2\le\log|D_n|\le A_*(100n)^2+E_n,
 \qquad E_n/n^2\longrightarrow0.
\]
Taking $n$ large enough that $E_n<n^2$ gives
\[
 \frac{1853339}{5000}n^2=370.6678n^2
 \le E_n<n^2,
\]
which is impossible. This proves Theorem~\ref{thm:main}, using the finite
certificates described above. Neither the argument nor its formalization
identifies which of the four values is irrational, proves irrationality
of an individual value, or establishes a three-value assertion.

The normalized separation is exactly
\[
 P_*-A_* =\frac{1853339}{50000000}=0.03706678.
\]
It compares asymptotic rates normalized by $(100n)^2$; it is not an error
bound for a numerical approximation to a zeta value or a determinant.
The proof needs an eventual cutoff, not its explicit numerical size.

\section{Formal verification and finite computations}\label{sec:verification}
\subsection{The theorem and its explicit hypotheses}
The final formal comparison takes a finite arithmetic grid $G$ with
predicate $G.\mathrm{Valid}$ and an analytic finite-knot certificate as
explicit hypotheses. All transitions from those hypotheses to
irrationality are checked in Lean~4 \cite{Lean4}. The two principal
explicit-hypothesis declarations, in namespace
\path{OddZeta.FourValue}, are
\begin{quote}\small
\path{not_all_rational_of_certificates} and
\path{irrational_one_of_four_of_certificates}.
\end{quote}
They occur in \path{FourValueFinalContradiction.lean}.
Their audited axiom dependencies are contained in
\texttt{propext}, \texttt{Classical.choice} and \texttt{Quot.sound};
neither depends on \texttt{sorryAx}.

The module \path{FourValueFinal.lean} supplies the finite certificates
through the isolated interfaces below. Its final declaration is
\begin{quote}\small
\path{OddZeta.FourValue.irrational_zeta7_or_zeta9_or_zeta11_or_zeta13}.
\end{quote}
It states the disjunction of Mathlib's \texttt{Irrational} predicates for
the real parts of the four classical \texttt{riemannZeta} values.
The proved lemmas \path{realZeta_cast} and \path{realZeta_eq_re} identify
these with $\sum_{m=1}^\infty m^{-k}$ for $k>1$.
Thus the conclusion concerns the usual real zeta values. The supplied
wrapper depends on \texttt{sorryAx}, precisely because the numerical
interfaces have not been replaced by Lean proof terms.

\subsection{The eight admitted finite families}
There are eight marked proof admissions in three project files, not eight
unproved asymptotic arguments. Each family is finite and has explicit data
and an independently runnable verifier.
\begin{center}\small
\begin{tabular}{p{.43\textwidth}cp{.40\textwidth}}
\toprule Interface file & Number & Finite content\\\midrule
\path{FourValueExternalTranscendental.lean}&2&Logarithm and arctangent interval enclosures.\\[5pt]
\path{FourValueExternalArithmetic.lean}&5&Rational identities and inequalities connecting the analytic grid to the knot bounds.\\[5pt]
\path{FourValueExternalGlobalArithmetic.lean}&1&Existence of the finite arithmetic grid satisfying its exact rational predicate.\\
\bottomrule
\end{tabular}
\end{center}
The source audit checks the marked locations and pinned statement hashes.
The log/arctangent checker uses outward-rounded Arb balls at 256-bit
precision and verifies containment in the claimed rational intervals
\cite{Arb}. The other two checkers use arbitrary-precision exact integer
and rational arithmetic. No optimizer output or floating-point sampling
is accepted as a substitute for these checks.

The completed reproduction run verified the following quantities.
\begin{center}\small
\begin{tabular}{p{.33\textwidth}p{.57\textwidth}}
\toprule Check & Verified coverage\\\midrule
Transcendental table & $5{,}700{,}022$ logarithm bounds and $1{,}200{,}000$ arctangent bounds.\\[4pt]
Analytic full grid & $50{,}000$ cells and $2{,}100{,}000$ knot checks, together with the exact argument, distance, cutoff and log-pair inequalities.\\[4pt]
Arithmetic grid & $171{,}504$ cells and $3{,}810{,}331$ reconstructed boxes, including zero-capacity boxes.\\
\bottomrule
\end{tabular}
\end{center}
All checks passed. The software checks establish the finite inputs
externally; they do not remove the admitted dependencies in Lean.
The proofs of contour convergence, rank, selected-minor transfer, valuation
assembly, boundary allowances, prime-sum limits and the final contradiction
are outside these finite admissions.

\subsection{Audit scope and dependencies}
The final axiom audit covers $1{,}873$ named declarations. Among the nine
new declarations in the final integration, only the three supplied-certificate
conclusions depend on \texttt{sorryAx}. The explicit-certificate comparison
has no such dependency. The prime-number-theorem source contains two
unrelated admitted declarations in its Wiener module; the audited
conditional global lower-bound theorem does not depend on them.

The build uses Lean 4.33.1, a pinned Mathlib revision \cite{Mathlib}, and
\texttt{PrimeNumberTheoremAnd} \cite{PNT}. The reproduction script completed
end to end, including fresh numerical checks and a fresh axiom audit.
It reused valid compiled artifacts and is not an empty-cache rebuild of
every historical source. An interrupted earlier run is recorded separately;
its partial log is not treated as successful verification. The final
archive was extracted independently and all 528 manifest entries matched.

The manuscript is an exposition of the construction and its verification
boundary. The exact formal statements, numerical predicates and data are
in the companion Round~29 package \cite{Project}. Formal checking verifies
those statements; editorial correspondence with the exposition and
independent mathematical review remain distinct responsibilities.

\appendix
\section{Exact polynomial-basis witness}\label{app:witness}
The following tables reproduce the finite witness used for the real
bound. Centers are indexed from 1. Every omitted weight is zero,
and every displayed weight numerator is divided by $10^8$. Row $j$
belongs to $s_j=C(2j+1)/80$, with $j$ indexed from 0. The first mass
vector at $s=0$ is zero; the terminal mass vector at $s=C$ uses row 39.
The radius is $1/50$, and the ray slope is $3/5$.
The corresponding data file is
\path{numerical/replay/round7_basis_witness.json} in the proof package.
\small
\begin{longtable}{rcc@{\hspace{12mm}}rcc}
\toprule $l$ & $\Re\xi_l$ & $\Im\xi_l$ & $l$ & $\Re\xi_l$ & $\Im\xi_l$\\ \midrule\endhead
1 & $\frac{103}{100}$ & $0$ & 24 & $\frac{3179}{2500}$ & $\frac{2563}{4000}$ \\[4pt]
2 & $\frac{2649}{2500}$ & $\frac{203}{4000}$ & 25 & $\frac{814}{625}$ & $\frac{177}{250}$ \\[4pt]
3 & $\frac{1287}{1250}$ & $\frac{203}{4000}$ & 26 & $\frac{3231}{2500}$ & $\frac{177}{250}$ \\[4pt]
4 & $\frac{2721}{2500}$ & $\frac{103}{1000}$ & 27 & $\frac{1653}{1250}$ & $\frac{3107}{4000}$ \\[4pt]
5 & $\frac{1323}{1250}$ & $\frac{103}{1000}$ & 28 & $\frac{3281}{2500}$ & $\frac{3107}{4000}$ \\[4pt]
6 & $\frac{2621}{2500}$ & $\frac{103}{1000}$ & 29 & $\frac{1677}{1250}$ & $\frac{847}{1000}$ \\[4pt]
7 & $\frac{679}{625}$ & $\frac{627}{4000}$ & 30 & $\frac{3329}{2500}$ & $\frac{847}{1000}$ \\[4pt]
8 & $\frac{2691}{2500}$ & $\frac{627}{4000}$ & 31 & $\frac{34}{25}$ & $\frac{147}{160}$ \\[4pt]
9 & $\frac{696}{625}$ & $\frac{53}{250}$ & 32 & $\frac{27}{20}$ & $\frac{147}{160}$ \\[4pt]
10 & $\frac{2759}{2500}$ & $\frac{53}{250}$ & 33 & $\frac{861}{625}$ & $\frac{124}{125}$ \\[4pt]
11 & $\frac{57}{50}$ & $\frac{43}{160}$ & 34 & $\frac{3419}{2500}$ & $\frac{124}{125}$ \\[4pt]
12 & $\frac{113}{100}$ & $\frac{43}{160}$ & 35 & $\frac{1743}{1250}$ & $\frac{4267}{4000}$ \\[4pt]
13 & $\frac{1457}{1250}$ & $\frac{327}{1000}$ & 36 & $\frac{3461}{2500}$ & $\frac{4267}{4000}$ \\[4pt]
14 & $\frac{2889}{2500}$ & $\frac{327}{1000}$ & 37 & $\frac{1763}{1250}$ & $\frac{1143}{1000}$ \\[4pt]
15 & $\frac{744}{625}$ & $\frac{1547}{4000}$ & 38 & $\frac{3501}{2500}$ & $\frac{1143}{1000}$ \\[4pt]
16 & $\frac{2951}{2500}$ & $\frac{1547}{4000}$ & 39 & $\frac{891}{625}$ & $\frac{4883}{4000}$ \\[4pt]
17 & $\frac{759}{625}$ & $\frac{56}{125}$ & 40 & $\frac{3539}{2500}$ & $\frac{4883}{4000}$ \\[4pt]
18 & $\frac{3011}{2500}$ & $\frac{56}{125}$ & 41 & $\frac{36}{25}$ & $\frac{13}{10}$ \\[4pt]
19 & $\frac{1547}{1250}$ & $\frac{2043}{4000}$ & 42 & $\frac{143}{100}$ & $\frac{13}{10}$ \\[4pt]
20 & $\frac{3069}{2500}$ & $\frac{2043}{4000}$ & 43 & $\frac{1817}{1250}$ & $\frac{5523}{4000}$ \\[4pt]
21 & $\frac{63}{50}$ & $\frac{23}{40}$ & 44 & $\frac{3609}{2500}$ & $\frac{5523}{4000}$ \\[4pt]
22 & $\frac{5}{4}$ & $\frac{23}{40}$ & 45 & $\frac{3641}{2500}$ & $\frac{1463}{1000}$ \\[4pt]
23 & $\frac{801}{625}$ & $\frac{2563}{4000}$ &  &  &  \\[4pt]
\bottomrule\end{longtable}
\begin{longtable}{r p{0.82\textwidth}}
\toprule Row $j$ & Nonzero pairs $(l,\text{weight numerator})$\\\midrule\endhead
0 & $(2,82721227), (6,17278773)$ \\[2pt]
1 & $(3,41767274), (5,51948099), (8,6284627)$ \\[2pt]
2 & $(3,36687301), (5,38270340), (7,6654700), (8,18387659)$ \\[2pt]
3 & $(2,931974), (3,31825824), (5,32215810), (7,24752571)$ \\[2pt]
 & $(10,10273821)$ \\[2pt]
4 & $(2,2032640), (3,28293195), (5,27867542), (7,22703973)$ \\[2pt]
 & $(9,18251540), (12,851110)$ \\[2pt]
5 & $(1,934855), (3,26741462), (5,25321036), (7,21741305)$ \\[2pt]
 & $(9,16685777), (12,8575565)$ \\[2pt]
6 & $(1,2092402), (3,23879064), (5,23056982), (7,20261790)$ \\[2pt]
 & $(9,16098475), (11,12938511), (14,1672776)$ \\[2pt]
7 & $(1,2866677), (3,21695570), (5,21295468), (7,18985445)$ \\[2pt]
 & $(9,16031508), (11,12340093), (14,6785239)$ \\[2pt]
8 & $(1,3674619), (2,2104956), (3,18108154), (5,19338360)$ \\[2pt]
 & $(7,17834262), (9,15366695), (11,10585185), (12,1514413)$ \\[2pt]
 & $(13,10033804), (16,1439552)$ \\[2pt]
9 & $(1,4251483), (2,3679060), (3,15338833), (5,17760838)$ \\[2pt]
 & $(7,16872551), (9,14802052), (11,12182820), (12,127587)$ \\[2pt]
 & $(13,9802262), (16,5182514)$ \\[2pt]
10 & $(1,4715553), (2,5388454), (3,12689033), (5,16299519)$ \\[2pt]
 & $(7,15994628), (9,14175381), (11,12121266), (13,5220122)$ \\[2pt]
 & $(14,4233932), (15,8118540), (18,1043572)$ \\[2pt]
11 & $(1,4890505), (2,5235437), (3,11800023), (5,15427635)$ \\[2pt]
 & $(7,15256356), (9,13609468), (11,12029010), (13,9162349)$ \\[2pt]
 & $(14,633783), (15,8009533), (18,3945901)$ \\[2pt]
12 & $(1,5239511), (2,7260414), (3,9246049), (5,14149161)$ \\[2pt]
 & $(7,14544455), (9,13145072), (11,11520299), (13,10027622)$ \\[2pt]
 & $(15,3617773), (16,3912547), (17,6713412), (20,623685)$ \\[2pt]
13 & $(1,5303690), (2,7004226), (3,8691797), (5,13501365)$ \\[2pt]
 & $(7,13948432), (9,12691567), (11,11238574), (13,9949950)$ \\[2pt]
 & $(15,6772031), (16,1227567), (17,6688005), (20,2982796)$ \\[2pt]
14 & $(1,5433372), (2,7752766), (3,7413498), (5,12676729)$ \\[2pt]
 & $(7,13393084), (9,12154456), (11,11215013), (13,8696173)$ \\[2pt]
 & $(14,702454), (15,8682393), (17,4193235), (18,1864377)$ \\[2pt]
 & $(19,5465769), (22,356681)$ \\[2pt]
15 & $(1,5441148), (2,7489856), (3,7025015), (5,12165625)$ \\[2pt]
 & $(7,12898415), (9,11774994), (11,10883762), (13,8017436)$ \\[2pt]
 & $(14,1220043), (15,8784428), (17,5768133), (18,543969)$ \\[2pt]
 & $(19,5789136), (22,2198040)$ \\[2pt]
16 & $(1,5509655), (2,8057769), (3,6043590), (5,11505147)$ \\[2pt]
 & $(7,12431750), (9,11475334), (11,10385073), (13,9471491)$ \\[2pt]
 & $(15,6239541), (16,1737798), (17,7502581), (19,4031745)$ \\[2pt]
 & $(20,964013), (21,4509494), (24,135019)$ \\[2pt]
17 & $(1,5486318), (2,7800284), (3,5761395), (5,11089155)$ \\[2pt]
 & $(7,12014643), (9,11129861), (11,10121567), (13,9291683)$ \\[2pt]
 & $(15,6134887), (16,1768982), (17,7559694), (19,4140727)$ \\[2pt]
 & $(20,1061521), (21,5049509), (24,1589774)$ \\[2pt]
18 & $(1,5519445), (2,8283731), (3,4951572), (5,10525182)$ \\[2pt]
 & $(7,11628255), (9,10700682), (11,10168785), (13,8771574)$ \\[2pt]
 & $(15,8337128), (17,5813402), (18,918350), (19,6423532)$ \\[2pt]
 & $(21,4320117), (23,3623664), (26,14581)$ \\[2pt]
19 & $(1,5477711), (2,8034975), (3,4742744), (5,10179873)$ \\[2pt]
 & $(7,11270047), (9,10399197), (11,9919006), (13,8597670)$ \\[2pt]
 & $(15,8222551), (17,5182555), (18,1476559), (19,6632417)$ \\[2pt]
 & $(21,3589416), (22,745034), (23,4396828), (26,1133417)$ \\[2pt]
20 & $(1,5478612), (2,8333468), (3,4157045), (5,9732126)$ \\[2pt]
 & $(7,10931495), (9,10195342), (11,9453527), (13,8714464)$ \\[2pt]
 & $(15,7866451), (17,7033626), (19,6110945), (21,5171034)$ \\[2pt]
 & $(23,4074387), (25,2747478)$ \\[2pt]
21 & $(1,5434827), (2,8217991), (3,3907660), (4,81110)$ \\[2pt]
 & $(5,9346406), (7,10599860), (9,9942337), (11,9141107)$ \\[2pt]
 & $(13,8712794), (15,7554028), (17,7143466), (19,3927838)$ \\[2pt]
 & $(20,1858987), (21,5827491), (23,2840469), (24,825409)$ \\[2pt]
 & $(25,3848818), (28,789402)$ \\[2pt]
22 & $(1,5385122), (2,8071441), (3,3726287), (4,1130925)$ \\[2pt]
 & $(5,8254558), (7,10104381), (9,9685718), (11,8912566)$ \\[2pt]
 & $(13,8583731), (15,7398836), (17,7134328), (19,5158107)$ \\[2pt]
 & $(20,703211), (21,5703957), (23,4149062), (25,3851093)$ \\[2pt]
 & $(27,2046677)$ \\[2pt]
23 & $(1,5359796), (2,8496160), (3,3140625), (4,2421511)$ \\[2pt]
 & $(5,6837689), (7,9570351), (9,9330800), (11,9008882)$ \\[2pt]
 & $(13,7936776), (15,7742518), (17,5246871), (18,1237670)$ \\[2pt]
 & $(19,6612926), (21,4318262), (22,520325), (23,5204132)$ \\[2pt]
 & $(25,2694599), (26,417925), (27,3374361), (30,527821)$ \\[2pt]
24 & $(1,5307703), (2,8282826), (3,3033861), (4,2360943)$ \\[2pt]
 & $(5,6655821), (7,9330328), (9,9112318), (11,8817801)$ \\[2pt]
 & $(13,7790190), (15,7625972), (17,5217375), (18,1205092)$ \\[2pt]
 & $(19,6579411), (21,4436785), (22,451807), (23,5279672)$ \\[2pt]
 & $(25,3363612), (27,3619235), (29,1113416), (30,415832)$ \\[2pt]
25 & $(1,5269950), (2,8429425), (3,2682352), (4,2490952)$ \\[2pt]
 & $(5,6245721), (7,9051447), (9,8983686), (11,8316360)$ \\[2pt]
 & $(13,8107070), (15,7033080), (17,6901357), (19,3626989)$ \\[2pt]
 & $(20,2022492), (21,6066180), (23,4121265), (25,4645702)$ \\[2pt]
 & $(27,2344848), (28,361643), (29,2945574), (32,353907)$ \\[2pt]
26 & $(1,5219043), (2,8307505), (3,2541640), (4,2473054)$ \\[2pt]
 & $(5,6037956), (7,8828739), (9,8783144), (11,8144317)$ \\[2pt]
 & $(13,7963697), (15,6920574), (17,6823504), (19,3533818)$ \\[2pt]
 & $(20,2069549), (21,6076042), (23,4155588), (25,4769970)$ \\[2pt]
 & $(27,2901825), (29,3212110), (32,1237925)$ \\[2pt]
27 & $(1,5178627), (2,8632525), (3,2099499), (4,4084683)$ \\[2pt]
 & $(5,4493029), (7,8277223), (9,8481472), (11,8256160)$ \\[2pt]
 & $(13,7354961), (15,7268515), (17,6306482), (19,6150365)$ \\[2pt]
 & $(21,5214648), (23,4986099), (25,3218230), (26,800472)$ \\[2pt]
 & $(27,3872514), (29,1194911), (30,1362787), (31,2532046)$ \\[2pt]
 & $(34,234752)$ \\[2pt]
28 & $(1,5177812), (2,8610902), (3,1861586), (4,4047782)$ \\[2pt]
 & $(5,4304217), (7,8078381), (9,8300920), (11,8088021)$ \\[2pt]
 & $(13,7229422), (15,7149842), (17,6233399), (19,6089184)$ \\[2pt]
 & $(21,5201055), (23,4991424), (25,3094191), (26,972433)$ \\[2pt]
 & $(27,4008111), (29,1266353), (30,1427948), (31,2889126)$ \\[2pt]
 & $(34,977891)$ \\[2pt]
29 & $(1,5169892), (2,8904105), (3,1438511), (4,4172288)$ \\[2pt]
 & $(5,3912430), (7,7848789), (9,8198859), (11,7694092)$ \\[2pt]
 & $(13,7375751), (15,6778547), (17,6399390), (19,5793159)$ \\[2pt]
 & $(21,5381328), (23,2454609), (24,2177412), (25,4764551)$ \\[2pt]
 & $(27,2928266), (28,397153), (29,3674810), (31,2159711)$ \\[2pt]
 & $(33,2224591), (36,151756)$ \\[2pt]
30 & $(1,5166191), (2,8803676), (3,1275923), (4,4099017)$ \\[2pt]
 & $(5,3796092), (7,7677629), (9,8032620), (11,7548618)$ \\[2pt]
 & $(13,7247870), (15,6674699), (17,6315770), (19,5735518)$ \\[2pt]
 & $(21,5346805), (23,2470862), (24,2156983), (25,4781846)$ \\[2pt]
 & $(27,2702879), (28,663830), (29,3843384), (31,2015417)$ \\[2pt]
 & $(32,220511), (33,2635668), (36,788192)$ \\[2pt]
31 & $(1,5153398), (2,9065568), (3,900404), (4,4211756)$ \\[2pt]
 & $(5,3438298), (7,7473363), (9,7865941), (11,7404077)$ \\[2pt]
 & $(13,7121802), (15,6571712), (17,6232564), (19,5677644)$ \\[2pt]
 & $(21,5319773), (23,4765570), (25,2675817), (26,1615538)$ \\[2pt]
 & $(27,4124016), (29,2280247), (30,834382), (31,3266224)$ \\[2pt]
 & $(33,1954935), (35,1958100), (38,88871)$ \\[2pt]
32 & $(1,5145966), (2,8965917), (3,759115), (4,4140984)$ \\[2pt]
 & $(5,3338389), (7,7318239), (9,7713596), (11,7269505)$ \\[2pt]
 & $(13,7001978), (15,6472380), (17,6150066), (19,5616956)$ \\[2pt]
 & $(21,5277856), (23,4747467), (25,2706152), (26,1590380)$ \\[2pt]
 & $(27,4148693), (29,1822856), (30,1327459), (31,3449187)$ \\[2pt]
 & $(33,1511136), (34,484079), (35,2400301), (38,641343)$ \\[2pt]
33 & $(1,5138747), (2,8872377), (3,623751), (4,4073940)$ \\[2pt]
 & $(5,3242973), (7,7170400), (9,7568129), (11,7140624)$ \\[2pt]
 & $(13,6886781), (15,6376346), (17,6069733), (19,5557189)$ \\[2pt]
 & $(21,5236328), (23,4734501), (25,4390587), (27,1596761)$ \\[2pt]
 & $(28,2184206), (29,3897304), (31,2664882), (32,42518)$ \\[2pt]
 & $(33,2979153), (35,1755833), (37,1749266), (40,47671)$ \\[2pt]
34 & $(1,5127640), (2,8779983), (3,497554), (4,4009124)$ \\[2pt]
 & $(5,3152160), (7,7028339), (9,7427820), (11,7015729)$ \\[2pt]
 & $(13,6774425), (15,6281685), (17,5989296), (19,5495356)$ \\[2pt]
 & $(21,5190082), (23,4708062), (25,4381904), (27,1619456)$ \\[2pt]
 & $(28,2175387), (29,3929451), (31,2380151), (32,372144)$ \\[2pt]
 & $(33,3136480), (35,1169410), (36,628204), (37,2195843)$ \\[2pt]
 & $(40,534315)$ \\[2pt]
35 & $(1,5109553), (2,9012433), (3,175826), (4,4456563)$ \\[2pt]
 & $(5,2593737), (7,6779302), (9,7283678), (11,6887305)$ \\[2pt]
 & $(13,6662450), (15,6187022), (17,5909949), (19,5433317)$ \\[2pt]
 & $(21,5144366), (23,4681947), (25,4383592), (27,3930718)$ \\[2pt]
 & $(29,2896996), (30,686799), (31,3266334), (33,1621603)$ \\[2pt]
 & $(34,1058602), (35,2586769), (37,991898), (38,675850)$ \\[2pt]
 & $(39,1563929), (42,19462)$ \\[2pt]
36 & $(1,5095500), (2,8918804), (3,67857), (4,4372576)$ \\[2pt]
 & $(5,2530746), (7,6654079), (9,7154351), (11,6771492)$ \\[2pt]
 & $(13,6557201), (15,6097200), (17,5832206), (19,5371680)$ \\[2pt]
 & $(21,5095831), (23,4650286), (25,4366564), (27,3932263)$ \\[2pt]
 & $(29,2944546), (30,659467), (31,3302942), (33,1672898)$ \\[2pt]
 & $(34,1073678), (35,2675038), (38,1749011), (39,2000550)$ \\[2pt]
 & $(42,453234)$ \\[2pt]
37 & $(1,5084036), (2,8770317), (4,4282346), (5,2497717)$ \\[2pt]
 & $(7,6536848), (9,7031835), (11,6661120), (13,6456499)$ \\[2pt]
 & $(15,6010912), (17,5757170), (19,5311771), (21,5048190)$ \\[2pt]
 & $(23,4618563), (25,4348815), (27,3934360), (29,3648497)$ \\[2pt]
 & $(31,1950894), (32,1240399), (33,3108591), (35,1298616)$ \\[2pt]
 & $(36,1037949), (37,2461204), (39,1463604), (41,1431257)$ \\[2pt]
 & $(44,8490)$ \\[2pt]
38 & $(1,5072352), (2,8521045), (4,4177583), (5,2513121)$ \\[2pt]
 & $(7,6427488), (9,6914955), (11,6554616), (13,6358543)$ \\[2pt]
 & $(15,5926224), (17,5682709), (19,5251258), (21,4998766)$ \\[2pt]
 & $(23,4583594), (25,4326009), (27,3926964), (29,3655079)$ \\[2pt]
 & $(31,1984387), (32,1231095), (33,3146498), (35,977102)$ \\[2pt]
 & $(36,1404513), (37,2616079), (39,685113), (40,799372)$ \\[2pt]
 & $(41,1859665), (44,405870)$ \\[2pt]
39 & $(1,5061914), (2,8277184), (4,4289237), (5,2375656)$ \\[2pt]
 & $(7,6278800), (9,6799795), (11,6449747), (13,6263201)$ \\[2pt]
 & $(15,5843582), (17,5610314), (19,5191927), (21,4950335)$ \\[2pt]
 & $(23,4548739), (25,4303213), (27,3919443), (29,3666736)$ \\[2pt]
 & $(31,3291645), (33,2127166), (34,865133), (35,2767260)$ \\[2pt]
 & $(37,951166), (38,1267720), (39,2196458), (41,1392256)$ \\[2pt]
 & $(43,1305784), (45,5589)$ \\[2pt]
\bottomrule\end{longtable}
\normalsize


\section{Source map and reproducibility}\label{app:reproduction}
\subsection{Principal formal modules}
Paths in this table are relative to the proof package's \texttt{lean/}
directory. The table identifies entry points rather than every dependency.
\begin{longtable}{p{.45\textwidth}p{.47\textwidth}}
\toprule Module & Role\\\midrule\endhead
\path{FourValueUniformNonvanishing.lean}&Uniform rank and nonzero-minor selection.\\[4pt]
\path{FourValueCorrectedZeta.lean}&Classical zeta identity in the corrected coordinates.\\[4pt]
\path{FourValueCorrectedAssembly.lean}&Entrywise and selected-determinant identification.\\[4pt]
\path{FourValueNumericalBridge.lean}&Analytic bound with one uniform $o(n^2)$ error.\\[4pt]
\path{FourValueFullGrid.lean}&Finite knot certificate supplied from the isolated numerical interfaces.\\[4pt]
\path{FourValueGlobalCosts.lean}&Finite dual valuation inequality for the actual corrected matrix.\\[4pt]
\path{FourValueArithmeticCertificate.lean}&Rational box and cell predicates.\\[4pt]
\path{FourValueGlobalPrimeSum.lean}&Product formula and finite weighted Chebyshev sum.\\[4pt]
\path{FourValueGlobalArithmeticRate.lean}&Limit and eventual arithmetic rate.\\[4pt]
\path{FourValueGlobalArithmeticLower.lean}&Nonzero selected rational minor under rationality.\\[4pt]
\path{FourValueFinalContradiction.lean}&Final comparison and real-series identification.\\[4pt]
\path{FourValueFinal.lean}&Supplied-certificate classical zeta conclusion.\\
\bottomrule
\end{longtable}

\subsection{Reproducing the proof-package checks}
Use the companion \path{Four_Value_Lean_Round29.zip}, not the superseded
Round~8 programs. The source package for this manuscript contains the TeX
files and the final report; the much larger proof package contains the
Lean sources, numerical tables and checkers. From the extracted proof
package root, first check the delivered bytes, then install the pinned
Python dependency and run:
\begin{verbatim}
sha256sum --check SHA256_MANIFEST.txt
python3 -m pip install python-flint==0.9.0
bash numerical/final29/reproduce.sh
\end{verbatim}
Lean must be installed using \path{lean/lean-toolchain}. The script obtains
the dependency cache, reruns all three numerical verifiers, builds the
final dependency chain and runs the axiom audit. It clears old success
records before starting and regenerates the combined numerical summary
from the new reports. A failed command stops the run.

The finite data can also be checked separately:
\begin{verbatim}
python3 numerical/external26/verify_values.py \
  numerical/external26/transcendental_values.txt \
  --report numerical/external26/verification_round29.json
python3 numerical/fullgrid27/verify_grid.py \
  --out numerical/fullgrid27/verification_round29.json
python3 numerical/global28/verify_arithmetic.py
\end{verbatim}
Use normal Python execution, not \texttt{python -O}, which disables
assertions. A packaged-evidence and source audit is available as
\begin{verbatim}
python3 numerical/final29/audit_final.py finish
\end{verbatim}
That last command inspects recorded evidence; it does not itself rerun
Lean or the numerical verifiers.

The pinned revisions are
\begin{quote}\small
Lean: \texttt{leanprover/lean4:v4.33.1}.\\
Mathlib: \path{0df444a360eaa60ab8c11dca51a86af692955474}.\\
PrimeNumberTheoremAnd:
\path{650d31264be65f4cd6e70c45d8b25d86d482a761}.
\end{quote}
The exact dependency manifest, data hashes, logs and admission markers are
included in the proof archive. Replacing the eight finite admissions by
Lean proofs would be a further kernel-certification task, not a missing
step in the conditional mathematical reduction.

\subsection{Reproducing this manuscript}
From the manuscript-source directory, run twice:
\begin{verbatim}
pdflatex -interaction=nonstopmode -halt-on-error \
  Four_Odd_Zeta_Manuscript.tex
\end{verbatim}
The source includes the section files, bibliography and exact witness
tables. Historical claims are supported by the references below;
\cite{Project} is the internal computational record and is not presented
as an independently published research article.
